%%
%% Berliner Hochschule für Technik --  Abschlussarbeit
%%
%% Kapitel 1
%%
%%

\chapter{Grundlagen und verwandte Arbeiten}

Dieses Kapitel ordnet die für diese Arbeit herangezogenen Quellen nach ihrer Bedeutung für die Fragestellung. Die Reihenfolge folgt einer argumentativen Relevanzordnung: Ausgangspunkt ist die gesundheitsökonomische Begründung durch die tabakbedingten Kosten in Deutschland \cite{effertz_viarisio_kosten_rauchens_2015}. Direkt danach steht die Laborstudie, aus der die zentrale Intervention abgeleitet wird: Tabibnia et al. untersuchten in einem fMRT-Paradigma, ob das Benennen rauchbezogener Reize das dadurch ausgelöste Rauchverlangen reduziert \cite{tabibnia_cuelabeling_2026}. Es folgen Leitlinien und epidemiologische Quellen zur Einordnung digitaler Tabakentwöhnung und der Zielgruppe \cite{who_tobacco_cessation_guideline_2024,olk_esa_tabak_2026,bmg_rauchen_2026}, verwandte Arbeiten zu Augmented Reality, Virtual Reality und Cue Exposure \cite{dougan_ar_smoking_cessation_2025,pericot_valverde_vr_smoking_cessation_2019}, Mess- und Interpretationsgrundlagen für subjektives Rauchverlangen \cite{addiction_research_center_qsu,weinfurt_clinically_meaningful_2019}, versorgungspolitische und finanzielle Rahmenbedingungen \cite{bmg_gesund_bleiben_praevention_2025,gba_tabakentwoehnung_2025,baecke_gkv_bonusprogramme_2026} sowie methodische Grundlagen für Klassifikation, Information Retrieval und Stichprobenplanung \cite{james_islr_2021,manning_ir_2008,cohen_power_2013,giner_sorolla_power_2024}. Diese Relevanzordnung macht deutlich, dass die Arbeit keine App-Entwicklung als Selbstzweck verfolgt, sondern die technisch umsetzbare, datensparsame und methodisch prüfbare Übertragung eines Laborbefunds in den Alltag.

\section{Gesundheitsökonomischer und versorgungspolitischer Kontext}

Die Motivation der Arbeit liegt in einem vermeidbaren Gesundheits- und Kostenproblem. Effertz und Viarisio bezifferten die direkten jährlichen Kosten des Rauchens in Deutschland auf 25,41 Milliarden Euro. Ihre Berechnung beruht auf krankheitsspezifischen Kostenmodellen und umfasst insbesondere Kosten für ambulante und stationäre Behandlung, Arzneimittel, Pflege und Rehabilitation. Einschließlich indirekter Kosten durch Produktivitätsverluste wurde die volkswirtschaftliche Last deutlich höher veranschlagt \cite{effertz_viarisio_kosten_rauchens_2015}. Für die vorliegende Arbeit ist dabei nicht die exakte Fortschreibung dieser Kosten entscheidend, sondern die gesundheitsökonomische Begründung: Schon kleine, niedrigschwellige Beiträge zur Unterstützung des Rauchstopps können relevant sein, wenn sie skalierbar sind, geringe Grenzkosten verursachen und für viele Betroffene ohne hohe Zugangshürden verfügbar sind.

Auch aktuelle gesundheitspolitische Quellen ordnen Rauchen weiterhin als wesentliches Präventionsziel ein. Das Bundesministerium für Gesundheit beschreibt Rauchen als das größte vermeidbare Gesundheitsrisiko in Deutschland und nennt die Reduktion des Tabakkonsums als wichtiges gesundheitspolitisches Ziel \cite{bmg_rauchen_2026}. Der Epidemiologische Suchtsurvey 2024 zeigt, dass Tabakkonsum nicht auf junge Zielgruppen beschränkt ist. In der Gesamtgruppe der 18- bis 64-Jährigen liegt der Anteil der Raucherinnen und Raucher bei 21,8~\%, in den Altersgruppen von 30 bis 59 Jahren bei ungefähr einem Fünftel bis knapp einem Viertel \cite{olk_esa_tabak_2026}. Diese Altersverteilung stützt die für die Studie vorgesehene Altersgruppe von 30 bis 65 Jahren: Sie deckt einen epidemiologisch relevanten Bereich ab und knüpft zugleich an den in der Laborstudie berichteten größeren Effekt bei älteren Teilnehmenden an.

Die WHO-Leitlinie zur Tabakentwöhnung bei Erwachsenen empfiehlt digitale Interventionen als mögliche Unterstützung für Personen, die ihren Tabakkonsum beenden möchten. Dabei wird die Evidenz differenziert bewertet: Textnachrichten werden mit höherer Sicherheit empfohlen als Smartphone-Apps und KI-basierte Interventionen, für die die Gewissheit der Evidenz geringer ist \cite{who_tobacco_cessation_guideline_2024}. Für diese Arbeit folgt daraus eine wichtige Einschränkung. Die zu entwickelnde App darf nicht als etablierte Therapie dargestellt werden. Sie ist ein Prototyp zur experimentellen Prüfung eines spezifischen Wirkmechanismus im Alltag. Die WHO-Empfehlung legitimiert digitale Zugangswege, ersetzt aber keine Wirksamkeitsprüfung der konkreten App-Funktion.

In Deutschland werden Prävention und Tabakentwöhnung zusätzlich durch Krankenkassen, Bonusprogramme und neue Leistungsansprüche gerahmt. Das Bundesministerium für Gesundheit beschreibt Prävention als Aufgabe, die sowohl individuelles Verhalten als auch Lebenswelten adressiert, und erwähnt digitale Präventionsangebote als Teil des Zugangs zu gesundheitsfördernden Maßnahmen \cite{bmg_gesund_bleiben_praevention_2025}. Der Gemeinsame Bundesausschuss hat Details für einen neuen Leistungsanspruch auf Arzneimittel zur Tabakentwöhnung bei schwerer Tabakabhängigkeit definiert und Anforderungen an begleitende Programme beschrieben \cite{gba_tabakentwoehnung_2025}. Bonusprogramme gesetzlicher Krankenkassen können gesundheitsförderliches Verhalten finanziell honorieren, sind aber uneinheitlich und freiwillig \cite{baecke_gkv_bonusprogramme_2026}. Die CueLens-Studie steht nicht in Konkurrenz zu diesen Angeboten. Sie untersucht einen eng begrenzten, nicht-pharmakologischen und appbasierten Mechanismus, der perspektivisch nur dann versorgungsrelevant wäre, wenn Wirksamkeit, Sicherheit, Datenschutz und Nutzbarkeit überzeugend nachgewiesen werden.

\section{Rauchverlangen, Auslöserreize und Cue-Labeling}

Rauchverlangen ist für diese Arbeit die zentrale Zielgröße. Gemeint ist nicht der langfristige Abstinenzerfolg, sondern die subjektiv berichtete Stärke des akuten Verlangens nach dem Rauchen in einer konkreten Situation. Solches Verlangen kann durch interne und externe Auslöserreize beeinflusst werden. Externe Auslöser sind beispielsweise Zigaretten, Feuerzeuge, Aschenbecher, Rauch, rauchende Personen oder soziale Situationen, die mit Rauchen verbunden sind. Für eine mobile App sind diese Reize besonders relevant, weil Smartphones Bilder darstellen und, bei Nutzung der Kamera, alltagsnahe visuelle Kontexte erfassen können.

Die Forschungsfrage der Masterarbeit knüpft direkt an die Studie von Tabibnia et al. an. Ausgangspunkt dieser Arbeit ist Affect Labeling, also das sprachliche Benennen emotional bedeutsamer Erfahrungen oder Reize. Tabibnia et al. übertragen diesen Ansatz auf rauchbezogene Reize und sprechen von Cue-Labeling. In der Cue-Labeling-Bedingung sollten Teilnehmende nicht den eigenen inneren Zustand benennen, also nicht etwa \enquote{Verlangen}, sondern sichtbare Merkmale des Reizes, etwa \enquote{smoke} oder \enquote{puff}. Diese Entscheidung ist für CueLens zentral, weil die App keine psychotherapeutische Selbstexploration erzwingen soll, sondern eine kurze, objektnahe Benennung eines wahrgenommenen Auslösers \cite{tabibnia_cuelabeling_2026}.

Die Laborstudie umfasste 50 Erwachsene, die täglich rauchten, zwischen 20 und 65 Jahre alt waren und vor der fMRT-Untersuchung für zwölf Stunden nicht rauchen sollten. Die Aufgabe bestand aus vier Bedingungen. In der neutralen Matching-Bedingung wurden neutrale Bilder einander zugeordnet. In der Cue-Matching-Bedingung wurden rauchbezogene Bilder anderen rauchbezogenen Bildern zugeordnet. In der Cue-Labeling-Bedingung wurde einem rauchbezogenen Bild eines von zwei rauchbezogenen Wörtern zugeordnet. In der Gender-Labeling-Bedingung wurde als Kontrollbedingung das Geschlecht einer abgebildeten Person benannt. Verwendet wurden unter anderem Bilder aus etablierten Bilddatenbanken und zugekaufte Bildmaterialien \cite{tabibnia_cuelabeling_2026}. Damit liefert die Studie nicht nur den psychologischen Wirkmechanismus, sondern auch ein technisches Grundmuster für die App: Reiz, Auswahlmenge, Antwort, anschließende Selbsteinschätzung.

Die Ergebnisse zeigen zunächst, dass die rauchbezogenen Bilder tatsächlich Rauchverlangen auslösten. Das berichtete Verlangen war in der Cue-Matching-Bedingung höher als in der neutralen Matching-Bedingung. Entscheidend für diese Arbeit ist der Vergleich zwischen Cue-Matching und Cue-Labeling: Das berichtete Verlangen war während Cue-Labeling niedriger als während Cue-Matching, der Gesamteffekt war jedoch klein (Hedges' \(g=-0{,}11\)). Zusätzlich zeigte sich geringere Aktivierung im Precuneus, einem Areal, das in der Studie selbst positiv mit dem berichteten Rauchverlangen assoziiert war. Alter moderierte den Effekt. Ab einem Johnson-Neyman-Grenzwert von 46,7 Jahren war das Verlangen während Cue-Labeling niedriger als während Cue-Matching und näherte sich der neutralen Bedingung an; für diese ältere Teilgruppe wurde ein Effekt von \(g=-0{,}29\) berichtet \cite{tabibnia_cuelabeling_2026}. Für die Zielgruppe dieser Masterarbeit ist der Altersbefund relevant, darf aber nicht überinterpretiert werden, weil die Laborstudie insgesamt klein war und Teilnehmende über 65 Jahre ausgeschlossen wurden.

Für die App-Studie ist Cue-Labeling deshalb als prüfbare Hypothese zu behandeln. Ein Laborbefund garantiert keine Wirkung im Alltag. Laborbedingungen erlauben kontrollierte Reizdarbietung, einheitliche Zeitfenster, standardisierte Antwortoptionen, fMRT-kompatible Abläufe und geringere Ablenkung. Im Alltag variieren Kontext, Motivation, Nikotinstatus, Stress, Umgebung und Smartphone-Nutzung. Außerdem können Bilder auf einem Display reale Rauchreize nur begrenzt abbilden. Die App muss daher den Laborvergleich möglichst sauber nachbilden: Cue-Labeling und Cue-Matching benötigen vergleichbare Reize, vergleichbare Bedienwege, eine kontrollierte Reihenfolge und eine unmittelbare Erfassung des Rauchverlangens. Zugleich muss sie alltagstauglicher sein als ein Laborparadigma, also kurz, fehlertolerant und ohne ausführliche Einweisung nutzbar.

Als Messbezug dient der Questionnaire on Smoking Urges \cite{addiction_research_center_qsu}. Für eine längere klinische Studie wäre ein validiertes mehrdimensionales Instrument vorzuziehen. Für die hier geplante App-Studie spricht jedoch die Machbarkeit für eine sehr kurze Selbstberichtsskala von 0 bis 100. Jede zusätzliche Frage erhöht die Bedienlast und kann im Alltag zu Abbrüchen führen. Die Vereinfachung ist methodisch vertretbar, solange die Interpretation begrenzt bleibt: Die Studie misst kurzfristige, subjektiv berichtete Unterschiede zwischen zwei Aufgabenbedingungen, nicht die Diagnose einer Tabakabhängigkeit und nicht den langfristigen Behandlungserfolg.

\section{Digitale und immersive Interventionen zur Tabakentwöhnung}

Digitale Interventionen werden in der Tabakentwöhnung vor allem eingesetzt, um Reichweite, zeitliche Verfügbarkeit und niedrige Zugangsschwellen zu erreichen. Smartphone-Apps können Erinnerungen, Selbstbeobachtung, Rückmeldungen, psychoedukative Inhalte und situationsbezogene Unterstützung verbinden. Die WHO-Leitlinie zeigt jedoch auch, dass die Evidenz nicht für alle digitalen Formate gleich stark ist \cite{who_tobacco_cessation_guideline_2024}. Daraus folgt für CueLens eine klare Abgrenzung: Die App ist kein umfassendes Entwöhnungsprogramm, sondern ein experimentelles Werkzeug, das eine einzelne, theoretisch begründete Aufgabenform im Alltag testet.

Verwandt sind Arbeiten zu Cue Exposure in virtuellen und erweiterten Umgebungen. Cue Exposure beruht auf der wiederholten Darbietung suchtbezogener Reize, ohne dass die gewohnte Konsumhandlung erfolgt. Dadurch soll die Reaktion auf den Reiz abgeschwächt werden. Pericot-Valverde et al. beschreiben Cue Exposure Treatment (CET) als kontrollierte und wiederholte Exposition gegenüber drogenbezogenen Reizen mit dem Ziel, cue-induziertes Verlangen zu extingieren \cite{pericot_valverde_vr_smoking_cessation_2019}. Die Relevanz für CueLens liegt in der gemeinsamen Annahme, dass visuelle Rauchreize messbares Verlangen auslösen und interventionell bearbeitet werden können. Zugleich unterscheidet sich CueLens wesentlich von CET: Die App soll Reize nicht passiv wiederholt darbieten, sondern eine kurze aktive Benennungs- beziehungsweise Zuordnungsaufgabe auslösen.

Die randomisierte klinische Studie von Pericot-Valverde et al. ist für die Einordnung besonders wichtig, weil sie die Grenzen eines plausiblen Reizexpositionsansatzes zeigt. In der Studie wurden 102 behandlungssuchende Raucherinnen und Raucher auf eine sechswöchige kognitiv-verhaltenstherapeutische Gruppenbehandlung (CBT, \(n=52\)) oder CBT plus fünf individuelle VR-basierte Cue-Exposure-Sitzungen (CBT+CET, \(n=50\)) randomisiert \cite{pericot_valverde_vr_smoking_cessation_2019}. Die virtuellen Umgebungen enthielten alltagsnahe Rauchkontexte wie Pub, Frühstück oder Mittagessen zu Hause, Kaffee im Café, Restaurant, Warten auf der Straße und Fernsehen am Abend. Teilnehmende konnten sich in den Umgebungen bewegen, mit rauchbezogenen Objekten interagieren und gaben während der Exposition ihr Rauchverlangen auf einer visuellen Analogskala von 0 bis 100 an \cite{pericot_valverde_vr_smoking_cessation_2019}.

Das Ergebnis ist ambivalent. CBT+CET reduzierte cue-induziertes Rauchverlangen über die VR-Expositionssitzungen hinweg, verbesserte aber weder Retention noch Abstinenzraten gegenüber CBT allein. Zusätzlich war die Rückfallrate im 12-Monats-Follow-up in der CBT+CET-Gruppe höher als in der CBT-Gruppe (Wald-\(\chi^2=4{,}796\), \(p=0{,}029\)) \cite{pericot_valverde_vr_smoking_cessation_2019}. Für CueLens ergibt sich daraus eine methodische und ethische Konsequenz: Die kurzfristige Reduktion von Verlangen ist ein relevanter, aber begrenzter Endpunkt. Das Auslösen oder Darbieten von Rauchreizen darf nicht als harmlos vorausgesetzt werden. Eine App-Studie muss die Belastung gering halten, Abbruchmöglichkeiten vorsehen, keine Abstinenzversprechen machen und die Ergebnisse auf kurzfristiges Verlangen begrenzen.

Noch näher an mobilen Alltagssituationen liegt die Arbeit von Dougan et al. zu Augmented Reality als Unterstützung der Rauchentwöhnung \cite{dougan_ar_smoking_cessation_2025}. Das zugehörige Studienprotokoll beschreibt den Vorteil, virtuelle Rauchreize in die natürliche Umgebung der Nutzenden einzublenden. Damit soll das bekannte Problem adressiert werden, dass in einem Labor gelernte Extinktionseffekte im Alltag nicht zuverlässig generalisieren. CueLens wählt bewusst einen technisch einfacheren Ansatz. Statt Rauchreize dreidimensional in den Raum zu projizieren, werden Bilder präsentiert oder durch ein Klassifikationsmodell aus Nutzendenbildern erkannt. Diese Entscheidung reduziert Entwicklungsaufwand, Geräteabhängigkeit und Fehlerrisiken. Sie passt zum Rahmen einer Masterarbeit, in der Robustheit, Datenschutz und Studiendurchführung höher zu gewichten sind als maximale Immersion.

Die verwandten Arbeiten zeigen damit eine doppelte Lücke. Einerseits gibt es digitale und immersive Interventionen gegen Rauchverlangen, andererseits ist die spezifische Übertragung von Cue-Labeling in den Alltag bislang nicht etabliert. Diese Masterarbeit positioniert sich genau in dieser Lücke: Sie verbindet eine verhaltensmedizinische Mikrowirkung mit einer mobilen, datensparsamen Android-Implementierung und einer kontrollierten Alltagsmessung.

\section{Technische Grundlagen der KI-gestützten Reizerkennung}

Der informatische Anteil der Arbeit liegt nicht nur in der App-Erstellung, sondern in der begründeten Auswahl technischer Datenstrukturen, Klassifikationsverfahren und Evaluationsmetriken. Aus technischer Sicht kann die App zwei Arten von Reizen verwenden. Erstens können vordefinierte Bilder eingesetzt werden, deren Kategorie und Label bereits bekannt sind. Zweitens können Nutzende eigene Bilder oder Kamerabilder verwenden. Dann muss die App erkennen, ob und welche rauchbezogenen Objekte oder Szenen enthalten sind. Diese zweite Variante erfordert ein Klassifikationsmodell und ist technisch anspruchsvoller.

Die Laborstudie zeigt ein für die App nutzbares Datenmodell: Ein Zielreiz wird mit einer kleinen Menge möglicher Antworten kombiniert. Für CueLens lässt sich dies als Relation zwischen Bild-ID, Reizkategorie, Antworttyp, Labeltext, Distraktor und Bedingung modellieren. Ein Eintrag kann beispielsweise die Kategorie \enquote{Rauch}, das Label \enquote{Rauch}, einen Distraktor wie \enquote{Packung} und die Bedingung \enquote{Cue-Labeling} enthalten. Für Cue-Matching wird dagegen ein weiterer Reiz derselben oder einer kontrolliert ähnlichen Kategorie angeboten. Eine solche Struktur ist technisch einfach, aber für die Studienqualität entscheidend: Nur wenn Bildmaterial, Kategorien und Antwortoptionen nachvollziehbar versioniert sind, kann später geprüft werden, ob Bedingungsunterschiede tatsächlich auf Labeling und nicht auf unterschiedlich stark auslösende Bilder zurückgehen.

Statistical Learning unterscheidet Klassifikationsprobleme von Regressionsproblemen: Bei der Klassifikation wird eine kategoriale Ausgabe vorhergesagt, etwa ob ein Bild eine Zigarette, Rauch, ein Feuerzeug oder keinen relevanten Reiz enthält \cite{james_islr_2021}. Für CueLens ist nicht allein die Genauigkeit entscheidend. Ein falsch positiver Treffer kann dazu führen, dass einer Person ein Rauchreiz oder Label angeboten wird, obwohl das Bild keinen relevanten Reiz enthält. Das kann irritieren und die Akzeptanz senken. Ein falsch negativer Treffer verhindert dagegen eine mögliche Intervention in einer relevanten Situation. Daraus folgt, dass Sensitivität, Spezifität, Präzision und Recall getrennt betrachtet werden müssen. Die F1-Kennzahl kann als harmonisches Mittel von Präzision und Recall nützlich sein, darf aber nicht die inhaltliche Bewertung der Fehlertypen ersetzen. Formal können die Bildentscheidungen in einer Konfusionsmatrix dargestellt werden. Für eine Reizklasse gilt dann: Präzision ist \(TP/(TP+FP)\), Recall beziehungsweise Sensitivität ist \(TP/(TP+FN)\), und \(F1=2\cdot(\mathrm{Precision}\cdot\mathrm{Recall})/(\mathrm{Precision}+\mathrm{Recall})\). Bei ungleichen Klassenhäufigkeiten ist zusätzlich zu prüfen, ob ein hoher Gesamtwert vor allem durch die häufige Klasse \enquote{kein Rauchreiz} entsteht.

Information-Retrieval-Grundlagen sind für die Labelauswahl relevant. Wenn ein Bild mehrere mögliche Reize enthält, muss die App passende Beschreibungen oder Vergleichsreize auswählen. Manning et al. beschreiben hierfür allgemeine Prinzipien wie Tokenisierung, Indexierung, Gewichtung und Rangordnung von Treffern \cite{manning_ir_2008}. Übertragen auf CueLens bedeutet dies: Labels und Vergleichsbilder sollten nicht beliebig präsentiert werden, sondern aus einer strukturierten Reiz- und Labelmenge stammen. Eine einfache Datenstruktur kann aus Reizkategorien, Synonymen, Bild-IDs, Labeltexten und Schwellenwerten bestehen. Für jedes erkannte Objekt wird ein Score gespeichert. Die App kann anschließend nur Labels anbieten, deren Score einen Mindestwert überschreitet. Dadurch wird die Auswahl nachvollziehbar, testbar und später austauschbar.

Datenschutz und Ressourcenanforderungen sprechen für eine lokale Verarbeitung auf dem Android-Endgerät. Bilder aus der Kamera oder aus der Galerie können besonders sensible Informationen enthalten, auch wenn die Studie selbst nur Rauchverlangen auswertet. Werden diese Bilder nicht übertragen und nicht gespeichert, sinkt das Datenschutzrisiko erheblich. Zugleich begrenzt On-Device-Klassifikation die Modellgröße, den Energieverbrauch und die Rechenzeit. Für die technische Evaluation sind daher neben Klassifikationsmetriken auch Speicherbedarf, Inferenzzeit, Modellgröße, Energieverbrauch und Android-Kompatibilität zu dokumentieren. Ein Modell, das auf einem Entwicklungsrechner hohe Genauigkeit erreicht, ist für diese Arbeit ungeeignet, wenn es auf realen Smartphones zu langsam, zu groß oder zu instabil ist.

\section{Messung, Stichprobenplanung und klinische Einordnung}

Die geplante App-Studie vergleicht Cue-Labeling und Cue-Matching innerhalb derselben Personen. Dieses gepaarte Design ist sinnvoll, weil interindividuelle Unterschiede im generellen Rauchverlangen, in der Abhängigkeitsschwere und in der Reaktivität auf Bilder teilweise kontrolliert werden. Für die statistische Auswertung ist deshalb nicht der Unterschied zwischen zwei unabhängigen Gruppen maßgeblich, sondern die Verteilung der paarweisen Differenzen. Die Laborstudie liefert hierfür zwei Orientierungswerte: In der Gesamtstichprobe war der Unterschied zwischen Cue-Matching und Cue-Labeling klein, in der älteren Teilgruppe ab etwa 47 Jahren größer, aber weiterhin moderat \cite{tabibnia_cuelabeling_2026}. Umgerechnet von der dortigen Skala von 1 bis 5 auf die in CueLens vorgesehene Skala von 0 bis 100 entspricht eine Differenz von 0,1 Skalenpunkten ungefähr 2,5 Punkten und eine Differenz von 0,2 Skalenpunkten ungefähr 5 Punkten. Diese Größenordnung ist für die Planung wichtig, weil sehr kleine Effekte im Alltag leicht durch Kontextschwankungen überdeckt werden können.

Die Stichprobenplanung orientiert sich an klassischer Power-Analyse nach Cohen \cite{cohen_power_2013}. Zugleich ist die neuere methodische Diskussion zu berücksichtigen. Giner-Sorolla et al. betonen, dass Power eine Eigenschaft eines konkreten statistischen Tests ist und dass Stichprobengrößen nicht allein nach Faustregeln bewertet werden sollten \cite{giner_sorolla_power_2024}. Für die vorliegende Arbeit ist das wichtig, weil ein alltagsnahes gepaartes Design mit wiederholten Berichten pro Person effizienter sein kann als ein reines Between-Subjects-Design. Die für die Pilotstudie angestrebten 68 verwertbaren Teilnehmenden sind daher als begründete Pilotplanung zu verstehen, nicht als Garantie für einen belastbaren Nachweis. Ausfälle, technische Probleme, unvollständige Selbstberichte und Heterogenität der Alltagssituationen müssen in der Interpretation berücksichtigt werden.

Eine statistisch signifikante Differenz ist außerdem nicht automatisch klinisch bedeutsam. Weinfurt unterscheidet zwischen einer Veränderung, die für Patientinnen und Patienten wahrnehmbar ist, und einer Veränderung, die im jeweiligen Kontext als wertvoll beurteilt wird \cite{weinfurt_clinically_meaningful_2019}. Diese Unterscheidung ist für CueLens zentral. Eine mittlere Reduktion um wenige Punkte auf einer 0-bis-100-Skala kann statistisch nachweisbar sein, aber für einzelne Nutzende kaum spürbar oder im Alltag wenig wertvoll. Umgekehrt kann eine kleine, wiederholt verfügbare Entlastung in kritischen Situationen subjektiv nützlich sein, wenn die App wenig Aufwand erzeugt und keine relevanten Risiken mit sich bringt. Die Arbeit sollte daher vorsichtig formulieren: Primärer Endpunkt ist die kurzfristige Veränderung des berichteten Rauchverlangens nach Reizbearbeitung; Aussagen zu Abstinenz, Rückfallprävention oder Behandlung der Tabakabhängigkeit wären erst in weiterführenden Studien zulässig.

\section{Konsequenzen für Konzeption und Studie}

Aus der Literatur ergeben sich fünf konkrete Anforderungen an die Masterarbeit. Erstens muss die gesundheitsökonomische Motivation klar von der experimentellen Zielgröße getrennt werden. Tabak verursacht hohe gesellschaftliche Kosten \cite{effertz_viarisio_kosten_rauchens_2015}, doch die App-Studie prüft zunächst nur eine kurzfristige Reduktion von Rauchverlangen. Zweitens muss der Laborbefund zu Cue-Labeling möglichst kontrolliert in die App übertragen werden. Cue-Labeling und Cue-Matching benötigen vergleichbare Reize, vergleichbare Bedienwege und eine unmittelbare Erhebung des Verlangens \cite{tabibnia_cuelabeling_2026}. Drittens muss die App die Darbietung von Rauchreizen ethisch vorsichtig handhaben, weil Cue-Exposure-Ansätze zwar Verlangen kurzfristig reduzieren können, aber nicht zwangsläufig Abstinenz verbessern und in der Studie von Pericot-Valverde et al. mit höheren Rückfallraten verbunden waren \cite{pericot_valverde_vr_smoking_cessation_2019}. Viertens muss die technische Umsetzung alltagstauglich sein. Ein lokal ausgeführtes Modell, begrenzte Bildspeicherung, kurze Interaktionszeiten und robuste Fehlerbehandlung sind für die Studienqualität ebenso wichtig wie die Klassifikationsgüte. Fünftens muss die Auswertung zwischen statistischer Signifikanz, wahrnehmbarer Veränderung und praktischem Nutzen unterscheiden \cite{weinfurt_clinically_meaningful_2019,giner_sorolla_power_2024}.

Damit ist der Rahmen für die folgenden Kapitel gesetzt. Die rechtlichen Rahmenbedingungen betreffen insbesondere Gesundheitsdaten, Einwilligung, Bildverarbeitung und den verantwortbaren Umgang mit potenziell craving-auslösenden Reizen. Die Evaluation geeigneter KI-Modelle muss technische Metriken mit Fehlertypen und Ressourcenanforderungen verbinden. Die Appentwicklung hat den experimentellen Vergleich so umzusetzen, dass die Nutzenden im Alltag möglichst wenig belastet werden. Die Studienplanung und Auswertung müssen schließlich offenlegen, welche Aussagen aus einem Prototyp mit begrenzter Fallzahl zulässig sind und welche nicht.
